\subsection{A degree-linear bound for refined-Swan zeros}

The following estimate passed an independent root proof audit.

\begin{proposition}[F13e: zeros on the original exceptional curves]
Retain the fixed pair $X_0,D_0$, matrix $B$, and cutoff vector $t$
of F13. Let $f:X_1\to X_0$ be a connected finite \'etale cover
of degree $d$, and let $\chi$ be a finite $p$-primary character
on $U_1$. Write
\[
 S_\chi=\sum_a s_aC'_a
\]
for its Swan divisor. If $C'_a$ lies above $C_i$ with degree $d_a$
and $s_a>0$, let $Z_a$ be the zero divisor of its nonzero refined
Swan section
\[
 \omega_a\in H^0\!\left(C'_a,
 \Omega^1_{X_1}(\log D_1)|_{C'_a}
                  \otimes\mathcal O_{X_1}(S_\chi)|_{C'_a}\right).
\]
More precisely, $Z_a$ is defined by saturating the image of
$\mathcal O_{C'_a}\xrightarrow{\omega_a}$ in this rank-two
vector bundle. Then
\[
 0\leq\deg Z_a\leq d_a(-B_{ii})(t_i-s_a),\qquad
 \sum_{a:s_a>0}\deg Z_a
       \leq d\sum_i(-B_{ii})t_i.
 \tag{F13.9}
\]
In particular the total number of zeros, counted with their
vanishing multiplicities on the component normalizations, is
bounded linearly in the cover degree. If $s_a=t_i$, that section
has no zero on $C'_a$.
\end{proposition}
\begin{proof}
The section exists and is nonzero by the exact refined-Swan
input used in F10: Krishna--Majumder, Corollary~9.7 and
Theorem~9.4, with the regular $F$-finite hypotheses explicitly
checked there \cite{KrishnaMajumder2026}. Its saturated image
is a line subbundle of the displayed rank-two vector bundle.
The map from $\mathcal O_{C'_a}$ to that line bundle has divisor
$Z_a$, so the degree of the line bundle is $\deg Z_a$.

The residue sequence (F10.4) on $C'_a$ has line subbundle and
line quotient of degrees $2g(C'_a)-2+\nu'_a$ and $0$.
Every line subbundle of its middle term has degree at most
$k'_a=d_a k_i$. Indeed, its map to the line quotient is either
nonzero, in which case its degree is at most zero, or vanishes,
in which case it maps into the line subbundle and has degree at
most $2g(C'_a)-2+\nu'_a$. Tensoring by
$\mathcal O_{X_1}(S_\chi)|_{C'_a}$ therefore gives
\[
                 \deg Z_a\leq d_a k_i+S_\chi\cdot C'_a.
 \tag{F13.10}
\]

The components over a fixed $C_i$ are disjoint: their union is
the finite \'etale inverse image of the smooth curve $C_i$.
The projection formula consequently gives
\[
 (C'_a)^2=d_aB_{ii},\qquad
 \sum_{b:C'_b\mapsto C_j}(C'_b\cdot C'_a)=d_aB_{ij}
                     \quad(j\ne i).
\]
By F13b, $0\leq s_b\leq t_j$ for every component over $C_j$.
All intersection numbers between distinct components are
nonnegative. It follows that
\[
 S_\chi\cdot C'_a
 \leq d_a\left(s_aB_{ii}+\sum_{j\ne i}t_jB_{ij}\right).
\]
The identity $(-B)t=k$ transforms the right side of (F13.10)
into $d_a(-B_{ii})(t_i-s_a)$. This proves the first inequality.
For each $i$, the degrees of all components above $C_i$ sum to
$d$. Dropping the nonnegative $s_a$ terms and summing proves
the second inequality. The final assertion follows because an
effective divisor of degree zero on a proper smooth curve over
an algebraically closed field is zero.
\end{proof}

\paragraph{Exact limitation.}
This estimate concerns the refined-Swan sections on the fixed
resolution and its finite \'etale covers. It does not bound the
number of blowups in a clean-ramification resolution, nor the
cohomology created at those zeros. Those implications require
an additional argument, so (F13.5) and (F13.6) remain open.
